% =====================================================================
% Chapitre 4 — Déploiement et gestion des versions (C2.2.4)
% =====================================================================

\competence{C2.2.4}{Historique des versions et dernière version fonctionnelle, fiable et viable}{Non}

Cette section montre que l'application est déployable de manière reproductible et que les versions sont tracées. La mise en production effective n'a pas encore eu lieu~: elle est planifiée pour août-septembre 2026. Le processus décrit ci-dessous a été préparé et testé en local ; il sera appliqué lors de la première mise en production.

\section{Système de gestion de versions}
Le projet utilise Git avec GitLab. La branche \texttt{main} a vocation à contenir la dernière version stable. Chaque version fera l'objet d'un tag Git (\texttt{v1.0.0}, \texttt{v1.0.1}, etc.) suivant le standard de versionnement sémantique (\texttt{major.minor.patch}).

\section{Processus de déploiement prévu en production}
\begin{enumerate}
    \item La branche \texttt{develop} est fusionnée dans \texttt{main} via une merge request dont le pipeline CI doit être au vert.
    \item Un tag de version est créé sur le commit de merge.
    \item Le pipeline CD se déclenche sur ce tag et construit les images Docker de l'API et du frontend.
    \item Les images sont poussées sur le registre Docker (GitLab Container Registry).
    \item Sur le serveur de production, \texttt{docker compose pull} récupère les nouvelles images et \texttt{docker compose up -d} redémarre les conteneurs.
    \item Prisma applique automatiquement les migrations de base de données en attente au démarrage du conteneur API.
\end{enumerate}

Ce processus a été préparé et testé en environnement local via Docker Compose, l'infrastructure de production prévue étant une transposition directe de cet environnement (voir C2.1.1). La mise en production réelle interviendra lors de la phase d'août-septembre 2026.

\section{Traçabilité des évolutions}
Les commits suivent la convention Conventional Commits (\texttt{feat:}, \texttt{fix:}, \texttt{chore:}, \texttt{docs:}, \texttt{test:}), ce qui rend l'historique lisible et exploitable pour générer un journal des changements. Le journal complet (\texttt{git log --oneline --decorate}) est fourni en annexe D.

\begin{table}[H]
    \centering
    \begin{tabular}{|>{\raggedright\arraybackslash}p{0.14\linewidth}|>{\raggedright\arraybackslash}p{0.16\linewidth}|>{\raggedright\arraybackslash}p{0.50\linewidth}|}
        \hline
        \textbf{Version} & \textbf{Date} & \textbf{Principales évolutions} \\
        \hline
        v1.0.0 & \textit{à renseigner} & Première version V1~: authentification, création de tournois, qualifications, saisie des résultats, brackets et classements. \\
        \hline
        v1.0.1 & \textit{à renseigner} & Corrections issues de la recette (voir C2.3.2). \\
        \hline
    \end{tabular}
\end{table}

\todo{Compléter le tableau des versions avec les dates et le contenu réels une fois les tags créés (production prévue août-septembre 2026).}

\section{Vérification de la stabilité à chaque déploiement}
À chaque déploiement, les scénarios de \emph{smoke test} du cahier de recettes seront exécutés manuellement~: connexion, création d'un tournoi simple, saisie d'un résultat et vérification du classement. Ce test de fumée prend moins de dix minutes et garantit que les fonctionnalités critiques sont opérationnelles.

\livrable{l'historique des différentes versions et la dernière version du logiciel fonctionnel, fiable et viable}
